\section{Conclusion}

This paper has described a workflow for building a high-precision corpus of public Slovene Bluesky posts and used that corpus for an initial look at the interactional and temporal structure of Slovene activity on the platform.

The pipeline combines relay-based tagged discovery (historical crawl and live firehose), DID-based backfill from individual Personal Data Servers, and post-level language filtering validated through manual annotation. It is fully protocol-native: no special API keys, no institutional data-access agreements, and no restriction to a single hosting server. For small-language communities where research infrastructure is limited, this openness is a practical advantage.

The corpus of 141,013 posts by 432 authors shows a sharp late-2024 expansion in recoverable activity, strong conversational use, and highly concentrated participation. Sparse hashtag use and link-sharing shaped largely by Slovene news media, GIF-based exchange, and platform-native tools further suggest a small public marked more by recurring interaction than by broad topic coordination.

This corpus should be treated as a snapshot of currently recoverable public posts from discovered Slovene-linked authors, not a complete census of all Slovene writing on Bluesky. Deleted posts are excluded by design, and within the pool of excluded untagged posts from known authors that meet the project's minimum-text rule, the false-negative rate is estimated at approximately 8.5\%. The main next step is not only technical enrichment, but a redistribution model that makes the corpus useful without ignoring the rights and expectations of the people whose posts make it possible.
